\section{Cómo leer el System Prompt de Codex}

\subsection{Método de lectura}

Para los usuarios que quieran entender a fondo las instrucciones internas de Codex, es posible:

\begin{enumerate}
    \item Localizar el archivo de sesión que guarda Codex (un archivo JSONL local)
    \item Extraer el contenido de cada mensaje de System, Developer y User
    \item El contenido en sí está en formato Markdown estándar, por lo que puede verse con una herramienta de renderizado de Markdown en línea
    \item Puede usarse una herramienta de traducción para leer el contenido en inglés capa por capa
\end{enumerate}

\begin{warningbox}{Ubicación del archivo de sesión}
El archivo de sesión de Codex se guarda en un directorio específico de la máquina local; la ruta exacta depende del sistema operativo y de la versión de Codex. Cada sesión corresponde a un archivo JSONL, en el que se registra en orden cronológico el historial completo de mensajes.
\end{warningbox}

\subsection{Resumen de la sección}

Al leer el archivo de sesión JSONL de Codex, es posible obtener el System Prompt completo, los mensajes de Developer y todo el contenido inyectado, lo cual es una vía eficaz para entender el diseño interno de un Modern Agent.

